% Part: lambda-calculus
% Chapter: lambda-definability
% Section: primitive-recursive-functions

\documentclass[../../../include/open-logic-section]{subfiles}

\begin{document}

\olfileid{lam}{ldf}{prf}
\olsection{توابع بازگشتی اولیه \usetoken{S}{lambda definable} هستند} 

به یاد آورید که توابع بازگشتی اولیه توابعی هستند که بتوان آن‌ها را از
توابع پایهٔ~$\Zero$، $\Succ$ و~$\Proj{n}{i}$ با ترکیب و بازگشت اولیه
تعریف کرد.

\begin{lem}
  \ollabel{lem:basic}
  توابع بازگشتی اولیهٔ پایه~$\Zero$، $\Succ$ و تصویرها~$\Proj{n}{i}$
  !!{lambda definable} هستند.
\end{lem}

\begin{proof}
ترم‌های زیر آن‌ها را~!!{lambda define}:
\begin{align*}
  \fn{Zero} & \ident \lambd[a][\lambd[fx][x]]\\
  \fn{Succ} & \ident \lambd[a][\lambd[fx][f (a f x)]] \\
  \fn{Proj}^n_i & \ident \lambd[x_0\dots x_{n-1}][x_i]
\end{align*}
\end{proof}

\begin{lem}
  \ollabel{lem:comp} فرض کنید تابع~$k$موضعی~$f$ و توابع~$n$موضعی
  $g_0, \dots, g_{k-1}$، به‌ترتیب با ترم‌های~$F$، $G_0$، \dots،
  $G_{k-1}$ !!{lambda definable} باشند، و~$h$ با ترکیب آن‌ها تعریف شده
  باشد. آنگاه~$h$ به‌وسیلهٔ~$H$ !!{lambda definable} است.
\end{lem}

\begin{proof}
  ترم زیر تابع~$h$ را~!!{lambda define}:
  \[
  H \ident \lambd[x_0 \dots x_{n-1}][F\, (G_0 x_0 \dots
    x_{n-1}) \dots (G_{k-1} x_0 \dots x_{n-1})]
  \]
  وارسی این واقعیت را به‌عنوان تمرین واگذار می‌کنیم.
\end{proof}

\begin{prob}
  برهان~\olref[lam][ldf][prf]{lem:comp} را با نشان‌دادن رابطهٔ زیر
  کامل کنید:
  $H\num{n_0}\dots\num{n_{n-1}} \red \num{h(n_0, \dots,
    n_{n-1})}$.
\end{prob}

توجه کنید که~\olref{lem:comp} مستلزم بازگشتی اولیه بودن~$f$ و~$g_0$،
\dots،~$g_{k-1}$ نبود؛ فقط لازم بود این توابع تام و
!!{lambda definable} باشند.

\begin{lem}\ollabel{lem:prim}
  فرض کنید~$f$ تابعی~$n$موضعی و~$g$ تابعی~$(n+2)$موضعی باشد، آن‌ها
  به‌ترتیب با ترم‌های~$F$ و~$G$ !!{lambda definable} باشند، و تابع~$h$
  با بازگشت اولیه از~$f$ و~$g$ تعریف شده باشد. آنگاه~$h$ نیز
  !!{lambda definable} است.
\end{lem}

\begin{proof}
  به یاد آورید که~$h$ چنین تعریف می‌شود:
  \begin{align*}
    h(x_1, \dots, x_n, 0) &= f(x_1, \dots, x_n)\\
    h(x_1, \dots, x_n, y+1) & = g(x_1, \dots, x_n, y,
      h(x_1, \dots, x_n, y)).
  \end{align*}
  به‌طور غیرصوری، تعریف بازگشتی اولیه، به‌روزرسانی حالت را که~$g$ تعیین
  می‌کند~$y$ بار تکرار می‌کند و از~$f(x_1, \dots, x_n)$ آغاز می‌شود.
  این امر یادآور تعریف عددنماهای چرچ است که آن‌ها نیز به‌صورت تکرارگر
  تعریف می‌شوند.

  برای سادگی، تعریف و برهان را در حالت~$n=1$ می‌آوریم: یک پارامتر~$x$
  به‌علاوهٔ متغیر بازگشت~$y$. تابع~$h$ با ترم زیر
  !!{lambda defined} است:
  \begin{align*}
    H \ident & \lambd[x][\lambd[y][\fn{Snd} (y
        D \tuple{\num{0}, F x})]]
    \intertext{که در آن}
    D \ident & \lambd[p][\tuple{\fn{Succ} (\fn{Fst}\, p), (G x
          (\fn{Fst}\, p) (\fn{Snd}\, p))}]
  \end{align*}
  حالت تکرار را به‌صورت زوجی نگه می‌داریم که مؤلفهٔ نخست آن مقدار کنونی
  $y$ و مؤلفهٔ دوم مقدار متناظر~$h$ است. در هر گام تکرار، با
  به‌روزرسانیِ تعیین‌شده به‌وسیلهٔ~$g$، زوجی از مقدار تازهٔ~$y$ و مقدار
  تازهٔ~$h$ می‌سازیم؛ پس از پایان تکرار، مؤلفهٔ دوم زوج را بازمی‌گردانیم
  که مقدار نهایی~$h$ است. اکنون ثابت می‌کنیم که این ترم واقعاً بازگشت
  اولیه را نمایش می‌دهد.

  می‌خواهیم ثابت کنیم برای هر~$n$ و~$m$،
  $H\,\num{n}\,\num{m} \red \num{h(n,m)}$. برای این کار نخست نشان
  می‌دهیم اگر~$D_n \ident \Subst{D}{\num{n}}{x}$، آنگاه
  $D_n^m \tuple{\num{0}, F\, \num{n}} \red
  \tuple{\num{m}, \num{h(n, m)}}$. با استقرا بر~$m$ پیش می‌رویم.

  اگر~$m=0$، باید نشان دهیم
  $D_n^0 \tuple{\num{0}, F\, \num{n}} \red
  \tuple{\num{0}, \num{h(n, 0)}}$. اما
  $D_n^0 \tuple{\num{0}, F\, \num{n}}$ همان
  $\tuple{\num{0}, F\, \num{n}}$ است. چون~$F$، تابع~$f$ را
  !!{lambda define}s، این ترم به~$\tuple{\num{0}, \num{f(n)}}$ کاهش
  می‌یابد؛ و چون~$f(n) = h(n, 0)$، این همان
  $\tuple{\num{0}, \num{h(n,0)}}$ است.

  اکنون فرض کنید~$D_n^m \tuple{\num{0}, F\, \num{n}} \red
  \tuple{\num{m}, \num{h(n, m)}}$. می‌خواهیم نشان دهیم
  $D_n^{m+1} \tuple{\num{0}, F\, \num{n}} \red
  \tuple{\num{m+1}, \num{h(n, m+1)}}$.
  \begin{align*}
    D_n^{m+1} \tuple{\num{0}, F\, \num{n}}
    & \ident D_n(D_n^{m} \tuple{\num{0}, F\, \num{n}})\\
    & \red D_n\,\tuple{\num{m}, \num{h(n, m)}} \text{\qquad (بنا بر فرض استقرا)}\\
    & \ident (\lambd[p][
      \tuple{
        \fn{Succ} (\fn{Fst}\, p), (G \, \num{n} (\fn{Fst}\, p) (\fn{Snd}\, p))
    }]) \tuple{\num{m}, \num{h(n,m)}}\\
    & \redone
    \tuple{\fn{Succ} (\fn{Fst}\, \tuple{\num{m}, \num{h(n,m)}}),\\
      & \qquad (G \, \num{n} (\fn{Fst}\, \tuple{\num{m}, \num{h(n,m)}}) (\fn{Snd}\, \tuple{\num{m}, \num{h(n,m)}}))}\\
    & \red 
    \tuple{\fn{Succ}\, \num{m}, (G \, \num{n} \,\num{m} \, \num{h(n,m)})}\\
    & \red \tuple{\num{m+1}, \num{g(n, m, h(n, m))}}
  \end{align*}
  چون~$g(n, m, h(n, m)) = h(n, m+1)$، کار تمام است.

  سرانجام، عبارت زیر را در نظر بگیرید:
  \begin{align*}
    H\,\num n\,\num m &\ident \lambd[x][\lambd[y][\fn{Snd} (y
        (\lambd[p].\tuple{\fn{Succ} (\fn{Fst}\, p), (G\, x\,
          (\fn{Fst}\, p)\, (\fn{Snd}\, p))}) \tuple{\num{0}, F x})]]\\
    & \qquad\,\num n\, \num m\\
    & \red \fn{Snd} (\num m \,
    \underbrace{(\lambd[p].\tuple{\fn{Succ} (\fn{Fst}\, p), (G \,\num n\,
      (\fn{Fst}\, p) (\fn{Snd}\, p))})}_{D_n} \tuple{\num{0}, F \num n})\\
    & \ident \fn{Snd} (\num m \, D_n\, \tuple{\num{0}, F \num n})\\
    & \red \fn{Snd} \,(D_n^m  \tuple{\num{0}, F \num n})\\
    & \red \fn{Snd} \,\tuple{\num{m}, \num{h(n, m)}} \\
    & \red \num{h(n,m)}. 
  \end{align*}  
\end{proof}


\begin{prop}
  هر تابع بازگشتی اولیه~!!{lambda definable} است.
\end{prop}

\begin{proof}
  بنا بر~\olref{lem:basic}، همهٔ توابع پایه~!!{lambda definable} هستند؛
  و بنا بر~\olref{lem:comp} و~\olref{lem:prim}، توابع
  !!{lambda definable} تحت ترکیب و بازگشت اولیه بسته‌اند.
\end{proof}

\end{document}
